% Section 4.5: The wedge and the tensor product | Part: 1 | Depends on: sec:tensors, sec:metric | Status: done | Original lines: 1176-1238
%!TEX root = ../../../main.tex
\subsection{The wedge and the tensor product}\label{sec:wedgetensor}

A $k$-tensor on $V$ is a $k$-linear map $V^k\to\R$ with no symmetry condition; these form $T^kV^*=V^*\otimes\cdots\otimes V^*$, of dimension
$n^k$. The tensor product multiplies them by juxtaposition, $(S\otimes T)(v_1,\dots,v_{k+l})=S(v_1,\dots,v_k)\,T(v_{k+1},\dots,v_{k+l})$;
it is associative but not commutative ($\dd x\otimes\dd y\neq\dd y\otimes\dd x$), and it obeys no relations.

\subhead{First description: the wedge is an antisymmetrised tensor product}
Define the antisymmetrisation (alternation) operator
\begin{equation}\label{eq:alt}
(\Alt T)(v_1,\dots,v_k)=\frac1{k!}\sum_{\sigma\in S_k}\sgn(\sigma)\,T\bigl(v_{\sigma(1)},\dots,v_{\sigma(k)}\bigr).
\end{equation}
It is a projection of $T^kV^*$ onto $\bigwedge^kV^*$ ($\Alt T=T$ exactly when $T$ is alternating). Comparing \eqref{eq:alt} for $T=\alpha\otimes\beta$
with \eqref{eq:wedge} gives
\begin{equation}\label{eq:altwedge}
\alpha\wedge\beta=\frac{(k+l)!}{k!\,l!}\,\Alt(\alpha\otimes\beta).
\end{equation}
\numcheck $k=l=1$: $\Alt(\dd x\otimes\dd y)=\tfrac12(\dd x\otimes\dd y-\dd y\otimes\dd x)$ and the prefactor is $2$, so
$\dd x\wedge\dd y=\dd x\otimes\dd y-\dd y\otimes\dd x$.

\begin{pitfall}[a warning on conventions]
% AUTHOR-NOTE: Kobayashi--Nomizu is named but has no bibliography entry (docs/AUTHOR_NOTES.md #9).
Some authors (\eg Kobayashi--Nomizu) drop the factor $(k+l)!/(k!\,l!)$, which gives $\dd x\wedge\dd y=\tfrac12(\dd x\otimes\dd y-\dd y\otimes\dd x)$ and
halves every formula above. The convention of this document for forms (that of Spivak and Lee~\cite{Spivak,Lee}) is the one for which
$(\dd x^1\wedge\cdots\wedge\dd x^n)(\partial_1,\dots,\partial_n)=1$, so that the volume form really gives volume $1$ to the unit cube.
The embedding of the multivector algebra $\bigwedge V\subset T(V)$ used in Section~\ref{sec:extquotient} uses the other normalisation
($v\wedge w\mapsto\frac12(v\otimes w-w\otimes v)$); the two differ by a factor $k!$ in degree $k$ and are isomorphic as algebras
(Section~\ref{sec:conventions}).
\end{pitfall}

\subhead{Second description: the tensor product with the single rule $\alpha\wedge\alpha=0$}
Let $T(V^*)=\bigoplus_kT^kV^*$ and let $\mathcal I$ be the two-sided ideal generated by all $\alpha\otimes\alpha$ with $\alpha\in V^*$. Then
\begin{equation}\label{eq:extquotient}
\textstyle\bigwedge(V^*)=\bigoplus_k\bigwedge^kV^*\ \cong\ T(V^*)/\mathcal I,\qquad \alpha\wedge\beta=[\alpha\otimes\beta].
\end{equation}
% REWRITTEN: identification stated once, in Remark rem:normalisation
This is the construction of Section~\ref{sec:extquotient} applied to $V^*$ (Remark~\ref{rem:normalisation}). The determinant is the prime example: $\bigwedge^nV^*$ is
one-dimensional, which is why, up to normalisation, there is only one alternating $n$-linear form on $\R^n$.

\subhead{Symmetric versus antisymmetric tensors}
Every $2$-tensor splits uniquely into a symmetric and an antisymmetric part, $V^*\otimes V^*=S^2V^*\oplus\bigwedge^2V^*$,
$T=\tfrac12(T+T^t)+\tfrac12(T-T^t)$. For example
\begin{equation}\label{eq:symantisym}
\dd x\otimes\dd y=\underbrace{\tfrac12(\dd x\otimes\dd y+\dd y\otimes\dd x)}_{\text{symmetric: }\dd x\cdot\dd y}+\underbrace{\tfrac12\,\dd x\wedge\dd y}_{\text{antisymmetric}}.
\end{equation}
\numcheck For $a=(1,2)$, $b=(3,4)$: $(\dd x\otimes\dd y)(a,b)=a^xb^y=4$; the symmetric part is $\frac12(a^xb^y+a^yb^x)=\frac12(4+6)=5$, the antisymmetric part is
$\frac12(4-6)=-1$, and $5+(-1)=4$ \checkmark. Also $(\dd x\wedge\dd y)(a,b)=4-6=-2=2\cdot(-1)$ \checkmark.

The two parts have different jobs. A \emph{metric} is a symmetric $2$-tensor, \eg $g=\dd x\otimes\dd x+\dd y\otimes\dd y$ (written $\dd x^2+\dd y^2$; this ``$\dd x^2$'' is a symmetric
product and has nothing to do with $\dd x\wedge\dd x=0$), and $g(v,v)\neq0$ measures length. A \emph{$2$-form} is an antisymmetric $2$-tensor, and $\omega(v,v)=0$ always; this is why $\iota_v\iota_v=0$ and why the
magnetic force does no work (Section~\ref{sec:phys-em}). For $k\ge3$ the space $T^kV^*$ contains further symmetry types, and $\bigwedge^k$ and $S^k$ are the two extreme ones.

\subhead{Duality with multivectors}
Dually, put $v\wedge w=v\otimes w-w\otimes v\in\bigwedge^2V$ (an oriented parallelogram, or bivector). The pairing of forms with multivectors is again a determinant,
\begin{equation}\label{eq:formmultipairing}
\langle\alpha^1\wedge\cdots\wedge\alpha^k,\ v_1\wedge\cdots\wedge v_k\rangle=\det\bigl[\alpha^i(v_j)\bigr]=(\alpha^1\wedge\cdots\wedge\alpha^k)(v_1,\dots,v_k),
\end{equation}
so contraction \eqref{eq:pairing} extends from $V^*\times V$ to $\bigwedge^kV^*\times\bigwedge^kV$. In this pairing the interior product is the transpose of wedging with $v$:
$\langle\iota_v\alpha,\xi\rangle=\langle\alpha,v\wedge\xi\rangle$ for $\xi\in\bigwedge^{k-1}V$ (this is just definition \eqref{eq:interior} below).

\begin{keyidea}[Why integrands are forms and not arbitrary tensors]
To assign a ``signed volume'' to a parallelepiped $(v_1,\dots,v_k)$ we need: (i) \emph{multilinearity}, so that volumes add when an edge is split;
(ii) \emph{alternation}, so that exchanging two edges flips the orientation and a degenerate (flat) parallelepiped has volume $0$: for a $2$-form $\omega(v,v)=0$ automatically, whereas a general $2$-tensor has $T(v,v)\neq0$ (a metric, for instance);
(iii) the \emph{Jacobian determinant} appearing under changes of variables, which \eqref{eq:detformula} produces by itself. Alternating multilinear maps are exactly $\bigwedge^kV^*\subset T^kV^*$, and the wedge product is the multiplication that keeps one inside this space.
\end{keyidea}
